\section{Cofinality}

\import{cardinal/infinite.tex}
\import{order/order-filter.tex}

For an ordinal, cofinality is defined by a set-bounded search over cardinal
domains.  On arbitrary sets these formulas give a deliberate total extension:
``cofinal in'' uses the subset relation literally, and no ordinal-theoretic
meaning is claimed off the ordinal domain.

\begin{definition}\label{cofinal_in}
    $\cofinalIn{C}{\alpha}$ iff $C\subseteq\alpha$ and
        for every $\beta\in\alpha$ there exists $\gamma\in C$ such that
        $\beta\subseteq\gamma$.
\end{definition}

\begin{definition}\label{cofinal_map}
    $\cofinalMap{f}{\delta}{\alpha}$ iff
        $f\in\funs{\delta}{\alpha}$ and
        $\cofinalIn{\ran{f}}{\alpha}$.
\end{definition}

\begin{theorem}\label{cofinal_map_iff_ordinal_cofinal_range}
    Let $\delta,\alpha$ be inhabited ordinals.
    Let $f\in\funs{\delta}{\alpha}$.
    Then $\cofinalMap{f}{\delta}{\alpha}$ iff
        $f$ has cofinal range from
            $\ordinalDirectedQuasiorder{\delta}$ to
            $\ordinalDirectedQuasiorder{\alpha}$.
\end{theorem}
\begin{proof}
    Let $J=\ordinalDirectedQuasiorder{\delta}$.
    Let $I=\ordinalDirectedQuasiorder{\alpha}$.
    $J$ is an exact directed quasiordered set by
        \auto{ordinal_directed_quasiorder_is_exact}.
    $I$ is an exact directed quasiordered set by
        \auto{ordinal_directed_quasiorder_is_exact}.
    $J$ is a quasiordered set by
        \obv{directed_quasiordered_set_is_quasiordered}.
    $I$ is a quasiordered set by
        \obv{directed_quasiordered_set_is_quasiordered}.
    $\carrier[J]=\delta$ by
        \auto{ordinal_directed_quasiorder_carrier}.
    $\carrier[I]=\alpha$ by
        \auto{ordinal_directed_quasiorder_carrier}.
    $\leq[I]=\subseteqrel{\alpha}$ by
        \auto{ordinal_directed_quasiorder_relation}.
    $f$ is a function by \obv{funs_is_function}.

    Show that if $\cofinalMap{f}{\delta}{\alpha}$, then
        $f$ has cofinal range from $J$ to $I$.
    \begin{subproof}
        Assume $\cofinalMap{f}{\delta}{\alpha}$.
        $\cofinalIn{\ran{f}}{\alpha}$ by \obv{cofinal_map}.
        Show that for every $\beta\in\carrier[I]$ there exists
            $j\in\carrier[J]$ such that
            $\beta\mathrel{\leq[I]}f(j)$.
        \begin{subproof}
            Fix $\beta\in\carrier[I]$.
            Hence $\beta\in\alpha$.
            Take $\gamma\in\ran{f}$ such that
                $\beta\subseteq\gamma$ by \obv{cofinal_in}.
            Take $j\in\dom{f}$ such that $f(j)=\gamma$ by
                \obv{fun_ran_iff}.
            Hence $j\in\delta$ by \obv{funs}.
            $f(j)\in\alpha$ by \obv{funs_type_apply}.
            $\beta\mathrel{\subseteqrel{\alpha}}f(j)$ by
                \obv{subseteqrel}.
            Hence $\beta\mathrel{\leq[I]}f(j)$.
        \end{subproof}
        $f\in\funs{\carrier[J]}{\carrier[I]}$.
        Follows by \obv{cofinal_range_between_quasiorders}.
    \end{subproof}

    Show that if $f$ has cofinal range from $J$ to $I$, then
        $\cofinalMap{f}{\delta}{\alpha}$.
    \begin{subproof}
        Assume $f$ has cofinal range from $J$ to $I$.
        $\ran{f}\subseteq\alpha$ by \obv{funs_ran}.
        Show that for every $\beta\in\alpha$ there exists
            $\gamma\in\ran{f}$ such that $\beta\subseteq\gamma$.
        \begin{subproof}
            Fix $\beta\in\alpha$.
            Hence $\beta\in\carrier[I]$.
            Take $j\in\carrier[J]$ such that
                $\beta\mathrel{\leq[I]}f(j)$ by
                \obv{cofinal_range_between_quasiorders}.
            Hence $j\in\delta$.
            $f(j)\in\ran{f}$ by \auto{fun_ran_iff,funs}.
            $\beta\mathrel{\subseteqrel{\alpha}}f(j)$.
            Hence $\beta\subseteq f(j)$ by \auto{subseteqrel}.
        \end{subproof}
        Hence $\cofinalIn{\ran{f}}{\alpha}$ by \obv{cofinal_in}.
        Follows by \obv{cofinal_map}.
    \end{subproof}
    Follows by \auto.
\end{proof}

\begin{corollary}\label{ordinal_at_top_cofinal_implies_cofinal_map}
    Let $\delta,\alpha$ be inhabited ordinals.
    Suppose $f$ is cofinal at the top from
        $\ordinalDirectedQuasiorder{\delta}$ to
        $\ordinalDirectedQuasiorder{\alpha}$.
    Then $\cofinalMap{f}{\delta}{\alpha}$.
\end{corollary}
\begin{proof}
    $\ordinalDirectedQuasiorder{\delta}$ is an exact directed
        quasiordered set by
        \auto{ordinal_directed_quasiorder_is_exact}.
    $\ordinalDirectedQuasiorder{\alpha}$ is an exact directed
        quasiordered set by
        \auto{ordinal_directed_quasiorder_is_exact}.
    $f\in\funs{\delta}{\alpha}$ by
        \auto{at_top_cofinal_type,
            ordinal_directed_quasiorder_carrier}.
    $f$ has cofinal range from
        $\ordinalDirectedQuasiorder{\delta}$ to
        $\ordinalDirectedQuasiorder{\alpha}$ by
        \obv{at_top_cofinal_implies_cofinal_range}.
    Follows by \obv{cofinal_map_iff_ordinal_cofinal_range}.
\end{proof}

\begin{theorem}\label{monotone_cofinal_map_is_ordinal_at_top_cofinal}
    Let $\delta,\alpha$ be inhabited ordinals.
    Let $f\in\funs{\delta}{\alpha}$.
    Suppose $\orderPreserving{f}{\delta}{\subseteqrel{\delta}}
        {\alpha}{\subseteqrel{\alpha}}$.
    Suppose $\cofinalMap{f}{\delta}{\alpha}$.
    Then $f$ is cofinal at the top from
        $\ordinalDirectedQuasiorder{\delta}$ to
        $\ordinalDirectedQuasiorder{\alpha}$.
\end{theorem}
\begin{proof}
    Let $J=\ordinalDirectedQuasiorder{\delta}$.
    Let $I=\ordinalDirectedQuasiorder{\alpha}$.
    $J$ is an exact directed quasiordered set by
        \auto{ordinal_directed_quasiorder_is_exact}.
    $I$ is an exact directed quasiordered set by
        \auto{ordinal_directed_quasiorder_is_exact}.
    $\carrier[J]=\delta$ by
        \auto{ordinal_directed_quasiorder_carrier}.
    $\carrier[I]=\alpha$ by
        \auto{ordinal_directed_quasiorder_carrier}.
    $\leq[J]=\subseteqrel{\delta}$ by
        \auto{ordinal_directed_quasiorder_relation}.
    $\leq[I]=\subseteqrel{\alpha}$ by
        \auto{ordinal_directed_quasiorder_relation}.
    Hence $\orderPreserving{f}{\carrier[J]}{\leq[J]}
        {\carrier[I]}{\leq[I]}$.
    $f$ has cofinal range from
        $\ordinalDirectedQuasiorder{\delta}$ to
        $\ordinalDirectedQuasiorder{\alpha}$ by
        \obv{cofinal_map_iff_ordinal_cofinal_range}.
    Hence $f$ has cofinal range from $J$ to $I$.
    Hence $f$ is cofinal at the top from $J$ to $I$ by
        \obv{monotone_cofinal_range_implies_at_top_cofinal}.
\end{proof}

\begin{definition}\label{cofinality_candidates}
    $\cofinalityCandidates{\alpha}
        =\{\kappa\in\suc{\card{\alpha}}\mid
            \text{$\kappa$ is a cardinal and there exists $f$ such that
            $\cofinalMap{f}{\kappa}{\alpha}$}\}$.
\end{definition}

\begin{definition}\label{cofinality}
    $\cofinality{\alpha}=\inters{\cofinalityCandidates{\alpha}}$.
\end{definition}

\begin{definition}\label{cofinal_upper_bounds}
    $\cofinalUpperBounds{f}{\xi}{\alpha}
        =\{\eta\in\alpha\mid
            \text{for every $j\in\suc{\xi}$ we have $f(j)\in\eta$}\}$.
\end{definition}

\begin{definition}\label{cofinal_envelope}
    $\cofinalEnvelope{f}{\lambda}{\alpha}
        =\{(\xi,\choice{\cofinalUpperBounds{f}{\xi}{\alpha}})
            \mid \xi\in\lambda\}$.
\end{definition}

\begin{proposition}\label{set_is_cofinal_in_itself}
    $\cofinalIn{A}{A}$.
\end{proposition}
\begin{proof}
    $A\subseteq A$ by \obv{subseteq_refl}.
    For every $a\in A$ there exists $b\in A$ such that $a\subseteq b$
        by \obv{subseteq_refl}.
    Follows by \obv{cofinal_in}.
\end{proof}

\begin{proposition}\label{surjection_is_cofinal_map}
    Let $f\in\surj{D}{A}$.
    Then $\cofinalMap{f}{D}{A}$.
\end{proposition}
\begin{proof}
    $f\in\funs{D}{A}$ by \obv{surj_to_fun}.
    $\ran{f}=A$ by \obv{ran_of_surj}.
    Hence $\cofinalIn{\ran{f}}{A}$
        by \obv{set_is_cofinal_in_itself}.
    Follows by \obv{cofinal_map}.
\end{proof}

\begin{proposition}\label{cofinality_candidates_inhabited}
    $\cofinalityCandidates{\alpha}$ is inhabited.
\end{proposition}
\begin{proof}
    Let $\kappa=\card{\alpha}$.
    $\kappa\in\suc{\kappa}$ by \obv{suc_intro_self}.
    $\kappa$ is a cardinal by \obv{card_is_cardinal}.
    $\kappa\approx\alpha$ by \auto{card_equinumerous}.
    Take $f$ such that $f\in\bijections{\kappa}{\alpha}$
        by \obv{equinum}.
    Hence $f\in\surj{\kappa}{\alpha}$ by \obv{bijections}.
    Hence $\cofinalMap{f}{\kappa}{\alpha}$
        by \obv{surjection_is_cofinal_map}.
    Hence $\kappa\in\cofinalityCandidates{\alpha}$
        by \obv{cofinality_candidates}.
\end{proof}

\begin{proposition}\label{cofinality_candidates_are_ordinals}
    Every element of $\cofinalityCandidates{\alpha}$ is an ordinal.
\end{proposition}
\begin{proof}
    Follows by
        \auto{cofinality_candidates,cardinal,initial_ordinal}.
\end{proof}

\begin{proposition}\label{cofinality_belongs_to_candidates}
    $\cofinality{\alpha}\in\cofinalityCandidates{\alpha}$.
\end{proposition}
\begin{proof}
    Follows by
        \auto{cofinality,inters_of_ordinals_elem,cofinality_candidates_inhabited,cofinality_candidates_are_ordinals}.
\end{proof}

\begin{theorem}\label{cofinality_is_cardinal}
    $\cofinality{\alpha}$ is a cardinal.
\end{theorem}
\begin{proof}
    Follows by
        \auto{cofinality_belongs_to_candidates,cofinality_candidates}.
\end{proof}

\begin{theorem}\label{cofinality_map_exists}
    There exists $f$ such that
        $\cofinalMap{f}{\cofinality{\alpha}}{\alpha}$.
\end{theorem}
\begin{proof}
    Follows by
        \auto{cofinality_belongs_to_candidates,cofinality_candidates}.
\end{proof}

\begin{proposition}\label{cofinality_subseteq_card}
    $\cofinality{\alpha}\subseteq\card{\alpha}$.
\end{proposition}
\begin{proof}
    $\cofinality{\alpha}\in\suc{\card{\alpha}}$
        by \auto{cofinality_belongs_to_candidates,cofinality_candidates}.
    Hence $\cofinality{\alpha}=\card{\alpha}$ or
        $\cofinality{\alpha}\in\card{\alpha}$ by \obv{suc_elim}.
    $\cofinality{\alpha}$ and $\card{\alpha}$ are ordinals
        by \auto{cofinality_is_cardinal,cardinal,initial_ordinal,card_is_ordinal}.
    Follows by \obv{subseteq_refl,ordinal_elem_implies_subseteq}.
\end{proof}

\begin{proposition}\label{cofinality_emptyset}
    $\cofinality{\emptyset}=\emptyset$.
\end{proposition}
\begin{proof}
    $\cofinality{\emptyset}\subseteq\card{\emptyset}$
        by \obv{cofinality_subseteq_card}.
    $\card{\emptyset}=\cardZero$ by \obv{card_of_cardinal_zero}.
    $\cardZero=\emptyset$ by \obv{cardinal_zero}.
    Hence $\cofinality{\emptyset}\subseteq\emptyset$.
    Follows by \obv{subseteq_emptyset_iff}.
\end{proof}

\begin{theorem}\label{cofinal_map_reindex}
    Suppose $\cofinalMap{f}{D}{\alpha}$ and $K\approx D$.
    Then there exists $g$ such that $\cofinalMap{g}{K}{\alpha}$.
\end{theorem}
\begin{proof}
    Take $b$ such that $b\in\bijections{K}{D}$ by \obv{equinum}.
    $b\in\funs{K}{D}$ by \obv{bijections_to_funs}.
    Hence $b$ is a function and $\dom{b}=K$
        by \obv{funs_is_function,funs}.
    $b\in\surj{K}{D}$ by \obv{bijections}.
    Hence $\ran{b}=D$ by \obv{ran_of_surj}.
    $f\in\funs{D}{\alpha}$ by \obv{cofinal_map}.
    Hence $f$ is a function and $\dom{f}=D$
        by \obv{funs_is_function,funs}.
    Let $g=f\circ b$.
    $g\in\funs{K}{\alpha}$ by \auto{funs_circ}.
    $\ran{g}=\ran{f}$ by \auto{ran_circ_exact}.
    $\cofinalIn{\ran{f}}{\alpha}$ by \obv{cofinal_map}.
    Hence $\cofinalIn{\ran{g}}{\alpha}$.
    Hence $\cofinalMap{g}{K}{\alpha}$ by \obv{cofinal_map}.
\end{proof}

\begin{theorem}\label{cofinality_le_domain_card}
    Suppose $\cofinalMap{f}{D}{\alpha}$.
    Then $\cofinality{\alpha}\subseteq\card{D}$.
\end{theorem}
\begin{proof}
    Let $c=\cofinality{\alpha}$.
    Let $\kappa=\card{\alpha}$.
    Let $\lambda=\card{D}$.
    $c\subseteq\kappa$ by \auto{cofinality_subseteq_card}.
    $\kappa$ and $\lambda$ are ordinals by \obv{card_is_ordinal}.
    $\lambda$ is a cardinal by \obv{card_is_cardinal}.
    $\lambda\approx D$ by \auto{card_equinumerous}.
    Take $g$ such that $\cofinalMap{g}{\lambda}{\alpha}$
        by \obv{cofinal_map_reindex}.
    \begin{byCase}
        \caseOf{$\kappa\subseteq\lambda$.}
            Follows by \obv{subseteq_transitive}.
        \caseOf{$\lnot(\kappa\subseteq\lambda)$.}
            Hence $\lambda\subseteq\kappa$
                by \obv{ordinal_subseteq_or_subseteq}.
            Hence $\lambda\in\kappa$ or $\lambda=\kappa$
                by \obv{ordinal_subseteq_implies_elem_or_eq}.
            Hence $\lambda\in\suc{\kappa}$ by \auto{suc_iff}.
            Hence $\lambda\in\cofinalityCandidates{\alpha}$
                by \obv{cofinality_candidates}.
            $c=\inters{\cofinalityCandidates{\alpha}}$
                by \obv{cofinality}.
            Follows by \auto{inters_subseteq_elem}.
    \end{byCase}
\end{proof}

\begin{proposition}\label{cofinality_subseteq_ordinal}
    Let $\alpha$ be an ordinal.
    Then $\cofinality{\alpha}\subseteq\alpha$.
\end{proposition}
\begin{proof}
    $\cofinality{\alpha}\subseteq\card{\alpha}$
        by \obv{cofinality_subseteq_card}.
    $\alpha\approx\alpha$ by \obv{equinum_refl}.
    Hence $\card{\alpha}\subseteq\alpha$
        by \obv{card_least_equinumerous_ordinal}.
    Follows by \obv{subseteq_transitive}.
\end{proof}

\begin{theorem}\label{cofinal_in_limit_iff_union}
    Let $\alpha$ be a limit ordinal.
    Suppose $C\subseteq\alpha$.
    Then $\cofinalIn{C}{\alpha}$ iff $\unions{C}=\alpha$.
\end{theorem}
\begin{proof}
    Show that if $\cofinalIn{C}{\alpha}$, then
        $\unions{C}=\alpha$.
    \begin{subproof}
        Assume $\cofinalIn{C}{\alpha}$.
        Show that every element of $\unions{C}$ is an element of $\alpha$.
        \begin{subproof}
            Fix $x\in\unions{C}$.
            Take $\gamma\in C$ such that $x\in\gamma$
                by \obv{unions_iff}.
            Hence $\gamma\in\alpha$ by \obv{subseteq}.
            Follows by \obv{ordinal_transitivity,limit_ordinal}.
        \end{subproof}
        Show that every element of $\alpha$ is an element of $\unions{C}$.
        \begin{subproof}
            Fix $\beta\in\alpha$.
            $\suc{\beta}\in\alpha$ by \obv{suc_elem_limitordinal}.
            Take $\gamma\in C$ such that $\suc{\beta}\subseteq\gamma$
                by \obv{cofinal_in}.
            $\beta\in\suc{\beta}$ by \obv{suc_intro_self}.
            Hence $\beta\in\gamma$ by \obv{subseteq}.
            Follows by \auto{unions_iff}.
        \end{subproof}
        Follows by set extensionality.
    \end{subproof}
    Show that if $\unions{C}=\alpha$, then
        $\cofinalIn{C}{\alpha}$.
    \begin{subproof}
        Assume $\unions{C}=\alpha$.
        Show that for every $\beta\in\alpha$ there exists $\gamma\in C$
            such that $\beta\subseteq\gamma$.
        \begin{subproof}
            Fix $\beta\in\alpha$.
            Hence $\beta\in\unions{C}$.
            Take $\gamma\in C$ such that $\beta\in\gamma$
                by \obv{unions_iff}.
            Hence $\gamma\in\alpha$ by \obv{subseteq}.
            $\beta$ and $\gamma$ are ordinals
                by \obv{limit_ordinal,elem_of_ordinal_is_ordinal}.
            Hence $\beta\subseteq\gamma$
                by \obv{ordinal_elem_implies_subseteq}.
        \end{subproof}
        Follows by \obv{cofinal_in}.
    \end{subproof}
    Follows by \auto.
\end{proof}

\begin{corollary}\label{cofinal_map_into_limit_iff_union_range}
    Let $\alpha$ be a limit ordinal.
    Let $f\in\funs{D}{\alpha}$.
    Then $\cofinalMap{f}{D}{\alpha}$ iff
        $\unions{\ran{f}}=\alpha$.
\end{corollary}
\begin{proof}
    $\ran{f}\subseteq\alpha$ by \obv{funs_ran}.
    Follows by \auto{cofinal_map,cofinal_in_limit_iff_union}.
\end{proof}

\begin{lemma}\label{small_domain_range_union_below}
    Let $\kappa$ be a limit ordinal.
    Let $\lambda=\cofinality{\kappa}$.
    Let $\beta\in\lambda$.
    Let $f\in\funs{\beta}{\kappa}$.
    Then $\unions{\ran{f}}\in\kappa$.
\end{lemma}
\begin{proof}
    Let $U=\unions{\ran{f}}$.
    $\ran{f}\subseteq\kappa$ by \obv{funs_ran}.
    $\kappa$ is an ordinal by \obv{limit_ordinal}.
    Hence every element of $\ran{f}$ is an ordinal
        by \obv{subseteq,elem_of_ordinal_is_ordinal}.
    Hence $U$ is an ordinal
        by \obv{union_of_ordinals_is_ordinal}.
    Show that every element of $U$ is an element of $\kappa$.
    \begin{subproof}
        Fix $x\in U$.
        Take $\gamma\in\ran{f}$ such that $x\in\gamma$
            by \obv{unions_iff}.
        Hence $\gamma\in\kappa$ by \obv{subseteq}.
        Follows by \obv{ordinal_transitivity}.
    \end{subproof}
    Hence $U\subseteq\kappa$ by \obv{subseteq}.
    Show that $U\neq\kappa$.
    \begin{subproof}
        Suppose not.
        Hence $\cofinalMap{f}{\beta}{\kappa}$
            by \auto{cofinal_map_into_limit_iff_union_range}.
        Hence $\cofinality{\kappa}\subseteq\card{\beta}$
            by \obv{cofinality_le_domain_card}.
        Hence $\lambda\subseteq\card{\beta}$.
        $\lambda$ is a cardinal by \obv{cofinality_is_cardinal}.
        Hence $\lambda$ is an ordinal by \auto{cardinal,initial_ordinal}.
        Hence $\beta$ is an ordinal
            by \obv{elem_of_ordinal_is_ordinal}.
        $\beta\approx\beta$ by \obv{equinum_refl}.
        Hence $\card{\beta}\subseteq\beta$
            by \obv{card_least_equinumerous_ordinal}.
        Hence $\lambda\subseteq\beta$
            by \obv{subseteq_transitive}.
        Hence $\beta\in\beta$ by \obv{subseteq}.
        Contradiction by \obv{in_irrefl}.
    \end{subproof}
    Follows by \auto{ordinal_subseteq_implies_elem_or_eq}.
\end{proof}

\begin{lemma}\label{finite_subset_of_limit_has_strict_upper_bound}
    Let $\alpha$ be a limit ordinal.
    Suppose $C$ is finite and $C\subseteq\alpha$.
    Then there exists $\eta\in\alpha$ such that
        for every $\gamma\in C$ we have $\gamma\in\eta$.
\end{lemma}
\begin{proof}
    Let $B=\{Y\in\pow{C}\mid
        \text{there exists $\beta\in\alpha$ such that
        for every $\gamma\in Y$ we have $\gamma\subseteq\beta$}\}$.
    $\emptyset\in\pow{C}$ by \obv{powerset_bottom}.
    $\emptyset\in\alpha$ by \obv{limit_ordinal}.
    Hence $\emptyset\in B$ by \auto{emptyset}.
    Show that for all $Y,y$ such that $Y\in B$ and $y\in C$ we have
        $\cons{y}{Y}\in B$.
    \begin{subproof}
        Fix $Y,y$.
        Assume $Y\in B$ and $y\in C$.
        Hence $Y\in\pow{C}$.
        Hence $Y\subseteq C$ by \obv{powerset_elim}.
        Hence $\cons{y}{Y}\subseteq C$
            by \obv{cons_subseteq_intro}.
        Hence $\cons{y}{Y}\in\pow{C}$ by \obv{powerset_intro}.
        Take $\beta\in\alpha$ such that
            for every $\gamma\in Y$ we have $\gamma\subseteq\beta$.
        Hence $y\in\alpha$ by \obv{subseteq}.
        Hence $\beta$ and $y$ are ordinals
            by \obv{elem_of_ordinal_is_ordinal}.
        Let $\delta=\beta\union y$.
        $\beta\subseteq y$ or $y\subseteq\beta$
            by \obv{ordinal_subseteq_or_subseteq}.
        Hence $\delta=y$ or $\delta=\beta$
            by \obv{union_absorb_subseteq_left,union_absorb_subseteq_right}.
        Hence $\delta\in\alpha$.
        For every $\gamma\in\cons{y}{Y}$ we have
            $\gamma\subseteq\delta$
            by \auto{cons_iff,union_upper_left,union_upper_right,subseteq_transitive}.
        Hence $\cons{y}{Y}\in B$.
    \end{subproof}
    Hence $C\in B$ by \obv{finite_induction}.
    Take $\beta\in\alpha$ such that
        for every $\gamma\in C$ we have $\gamma\subseteq\beta$.
    Let $\eta=\suc{\beta}$.
    $\eta\in\alpha$ by \auto{suc_elem_limitordinal}.
    $\beta$ is an ordinal
        by \obv{limit_ordinal,elem_of_ordinal_is_ordinal}.
    For every $\gamma\in C$ we have $\gamma$ is an ordinal
        by \obv{subseteq,limit_ordinal,elem_of_ordinal_is_ordinal}.
    Hence for every $\gamma\in C$ we have
        $\gamma\in\eta$
        by \auto{ordinal_subseteq_implies_elem_or_eq,suc_iff}.
\end{proof}

\begin{theorem}\label{cofinality_of_limit_is_infinite_cardinal}
    Let $\alpha$ be a limit ordinal.
    Then $\cofinality{\alpha}$ is an infinite cardinal.
\end{theorem}
\begin{proof}
    Let $\lambda=\cofinality{\alpha}$.
    $\lambda$ is a cardinal by \obv{cofinality_is_cardinal}.
    Show that $\lambda$ is infinite.
    \begin{subproof}
        Suppose not.
        Hence $\lambda$ is finite.
        Take $f$ such that $\cofinalMap{f}{\lambda}{\alpha}$
            by \auto{cofinality_map_exists}.
        Hence $f\in\funs{\lambda}{\alpha}$ and
            $\cofinalIn{\ran{f}}{\alpha}$ by \obv{cofinal_map}.
        Hence $f$ is a function and $\dom{f}=\lambda$
            by \obv{funs_is_function,funs}.
        Hence $\img{f}{\lambda}$ is finite
            by \obv{finite_function_image,subseteq_refl}.
        $\img{f}{\lambda}=\ran{f}$ by \auto{img_dom}.
        Hence $\ran{f}$ is finite.
        $\ran{f}\subseteq\alpha$ by \obv{funs_ran}.
        Take $\eta\in\alpha$ such that
            for every $\gamma\in\ran{f}$ we have $\gamma\in\eta$
            by \obv{finite_subset_of_limit_has_strict_upper_bound}.
        Take $\gamma\in\ran{f}$ such that $\eta\subseteq\gamma$
            by \obv{cofinal_in}.
        Hence $\gamma\in\eta$.
        Hence $\gamma\in\gamma$ by \obv{subseteq}.
        Contradiction by \obv{in_irrefl}.
    \end{subproof}
\end{proof}

\begin{corollary}\label{cofinality_of_limit_is_limit_ordinal}
    Let $\alpha$ be a limit ordinal.
    Then $\cofinality{\alpha}$ is a limit ordinal.
\end{corollary}
\begin{proof}
    Follows by
        \obv{cofinality_of_limit_is_infinite_cardinal,infinite_cardinal_is_limit_ordinal}.
\end{proof}

\begin{lemma}\label{cofinal_upper_bounds_inhabited}
    Let $\alpha$ be a limit ordinal.
    Let $\lambda=\cofinality{\alpha}$.
    Suppose $\cofinalMap{f}{\lambda}{\alpha}$.
    Let $\xi\in\lambda$.
    Then $\cofinalUpperBounds{f}{\xi}{\alpha}$ is inhabited.
\end{lemma}
\begin{proof}
    $\lambda$ is a limit ordinal
        by \obv{cofinality_of_limit_is_limit_ordinal}.
    Hence $\suc{\xi}\in\lambda$ by \obv{suc_elem_limitordinal}.
    $\lambda$ and $\suc{\xi}$ are ordinals
        by \obv{limit_ordinal,elem_of_ordinal_is_ordinal}.
    Hence $\suc{\xi}\subseteq\lambda$
        by \obv{ordinal_elem_implies_subseteq}.
    $f\in\funs{\lambda}{\alpha}$ by \obv{cofinal_map}.
    Hence $f$ is a function and $\dom{f}=\lambda$
        by \obv{funs_is_function,funs}.
    Let $r=\restrl{f}{\suc{\xi}}$.
    $r\in\funs{\suc{\xi}}{\alpha}$ by \auto{funs_restrl}.
    Show that not $\cofinalMap{r}{\suc{\xi}}{\alpha}$.
    \begin{subproof}
        Suppose $\cofinalMap{r}{\suc{\xi}}{\alpha}$.
        Hence $\cofinality{\alpha}\subseteq\card{\suc{\xi}}$
            by \obv{cofinality_le_domain_card}.
        Hence $\lambda\subseteq\card{\suc{\xi}}$.
        $\suc{\xi}\approx\suc{\xi}$ by \obv{equinum_refl}.
        Hence $\card{\suc{\xi}}\subseteq\suc{\xi}$
            by \obv{card_least_equinumerous_ordinal}.
        Hence $\lambda\subseteq\suc{\xi}$
            by \obv{subseteq_transitive}.
        Hence $\suc{\xi}\in\suc{\xi}$ by \obv{subseteq}.
        Contradiction by \obv{in_irrefl}.
    \end{subproof}
    Hence not $\cofinalIn{\ran{r}}{\alpha}$
        by \auto{cofinal_map}.
    $\ran{r}\subseteq\alpha$ by \obv{funs_ran}.
    Take $\eta\in\alpha$ such that for every $\gamma\in\ran{r}$
        we have $\lnot(\eta\subseteq\gamma)$ by \auto{cofinal_in}.
    Show that for every $j\in\suc{\xi}$ we have $f(j)\in\eta$.
    \begin{subproof}
        Fix $j\in\suc{\xi}$.
        $r(j)=f(j)$ by \auto{restrl_of_function_apply}.
        $r$ is a function and $\dom{r}=\suc{\xi}$
            by \obv{funs_is_function,funs}.
        Hence $r(j)\in\ran{r}$ by \auto{fun_ran_iff}.
        Hence $\lnot(\eta\subseteq f(j))$.
        $f(j)\in\alpha$ by \obv{funs_type_apply}.
        Hence $\eta$ and $f(j)$ are ordinals
            by \obv{elem_of_ordinal_is_ordinal}.
        $f(j)\subseteq\eta$ by \obv{ordinal_subseteq_or_subseteq}.
        $f(j)\neq\eta$.
        Follows by \obv{ordinal_subseteq_implies_elem_or_eq}.
    \end{subproof}
    Hence $\eta\in\cofinalUpperBounds{f}{\xi}{\alpha}$
        by \obv{cofinal_upper_bounds}.
\end{proof}

\begin{proposition}\label{cofinal_envelope_type}
    Let $\alpha$ be a limit ordinal.
    Let $\lambda=\cofinality{\alpha}$.
    Suppose $\cofinalMap{f}{\lambda}{\alpha}$.
    Then $\cofinalEnvelope{f}{\lambda}{\alpha}
        \in\funs{\lambda}{\alpha}$.
\end{proposition}
\begin{proof}
    Let $h=\cofinalEnvelope{f}{\lambda}{\alpha}$.
    $h$ is a relation by \auto{cofinal_envelope,relation}.
    $h$ is right-unique
        by \auto{cofinal_envelope,rightunique,pair_eq_iff}.
    $\dom{h}=\lambda$
        by \auto{cofinal_envelope,setext,dom_iff,pair_eq_iff}.
    For every $\xi\in\lambda$ we have
        $\cofinalUpperBounds{f}{\xi}{\alpha}$ is inhabited
        by \obv{cofinal_upper_bounds_inhabited}.
    Hence for every $\xi\in\lambda$ we have
        $\choice{\cofinalUpperBounds{f}{\xi}{\alpha}}
            \in\cofinalUpperBounds{f}{\xi}{\alpha}$
        by \obv{choice_in}.
    Hence for every $\xi\in\lambda$ we have
        $\choice{\cofinalUpperBounds{f}{\xi}{\alpha}}\in\alpha$
        by \auto{cofinal_upper_bounds}.
    For every $\xi\in\lambda$ we have
        $h(\xi)=\choice{\cofinalUpperBounds{f}{\xi}{\alpha}}$
        by \obv{cofinal_envelope,function_apply_intro}.
    Hence $h$ is a function to $\alpha$ by \auto.
    Follows by \auto{funs_intro}.
\end{proof}

\begin{proposition}\label{cofinal_envelope_application}
    Let $\alpha$ be a limit ordinal.
    Let $\lambda=\cofinality{\alpha}$.
    Suppose $\cofinalMap{f}{\lambda}{\alpha}$.
    Let $\xi\in\lambda$.
    Then $(\cofinalEnvelope{f}{\lambda}{\alpha})(\xi)
        =\choice{\cofinalUpperBounds{f}{\xi}{\alpha}}$.
\end{proposition}
\begin{proof}
    $(\xi,\choice{\cofinalUpperBounds{f}{\xi}{\alpha}})
        \in\cofinalEnvelope{f}{\lambda}{\alpha}$
        by \obv{cofinal_envelope}.
    $\cofinalEnvelope{f}{\lambda}{\alpha}$ is a function
        by \auto{cofinal_envelope_type,funs_is_function}.
    Follows by \obv{function_apply_intro}.
\end{proof}

\begin{proposition}\label{cofinal_envelope_dominates_prefix}
    Let $\alpha$ be a limit ordinal.
    Let $\lambda=\cofinality{\alpha}$.
    Suppose $\cofinalMap{f}{\lambda}{\alpha}$.
    Let $\xi\in\lambda$.
    Let $j\in\suc{\xi}$.
    Then $f(j)\in(\cofinalEnvelope{f}{\lambda}{\alpha})(\xi)$.
\end{proposition}
\begin{proof}
    $\cofinalUpperBounds{f}{\xi}{\alpha}$ is inhabited
        by \obv{cofinal_upper_bounds_inhabited}.
    Hence $\choice{\cofinalUpperBounds{f}{\xi}{\alpha}}
        \in\cofinalUpperBounds{f}{\xi}{\alpha}$ by \obv{choice_in}.
    Hence $f(j)\in\choice{\cofinalUpperBounds{f}{\xi}{\alpha}}$
        by \obv{cofinal_upper_bounds}.
    Follows by \obv{cofinal_envelope_application}.
\end{proof}

\begin{theorem}\label{cofinality_of_cofinality}
    Let $\alpha$ be a limit ordinal.
    Then $\cofinality{\cofinality{\alpha}}=\cofinality{\alpha}$.
\end{theorem}
\begin{proof}
    Let $\lambda=\cofinality{\alpha}$.
    Let $\mu=\cofinality{\lambda}$.
    $\lambda$ is a limit ordinal
        by \obv{cofinality_of_limit_is_limit_ordinal}.
    Take $f$ such that $\cofinalMap{f}{\lambda}{\alpha}$
        by \auto{cofinality_map_exists}.
    Take $p$ such that $\cofinalMap{p}{\mu}{\lambda}$
        by \auto{cofinality_map_exists}.
    Hence $f\in\funs{\lambda}{\alpha}$ and
        $\cofinalIn{\ran{f}}{\alpha}$ by \obv{cofinal_map}.
    Hence $p\in\funs{\mu}{\lambda}$ and
        $\cofinalIn{\ran{p}}{\lambda}$ by \obv{cofinal_map}.
    Let $h=\cofinalEnvelope{f}{\lambda}{\alpha}$.
    $h\in\funs{\lambda}{\alpha}$
        by \auto{cofinal_envelope_type}.
    Hence $h$ is a function and $\dom{h}=\lambda$
        by \obv{funs_is_function,funs}.
    $p$ is a function and $\dom{p}=\mu$
        by \obv{funs_is_function,funs}.
    $\ran{p}\subseteq\lambda$ by \obv{funs_ran}.
    Hence $h$ is composable with $p$.
    Let $g=h\circ p$.
    $g\in\funs{\mu}{\alpha}$ by \auto{funs_circ}.
    Show that $\cofinalIn{\ran{g}}{\alpha}$.
    \begin{subproof}
        $\ran{g}\subseteq\alpha$ by \obv{funs_ran}.
        Show that for every $\beta\in\alpha$ there exists
            $\rho\in\ran{g}$ such that $\beta\subseteq\rho$.
        \begin{subproof}
            Fix $\beta\in\alpha$.
            Take $\gamma\in\ran{f}$ such that $\beta\subseteq\gamma$
                by \obv{cofinal_in}.
            $f$ is a function and $\dom{f}=\lambda$
                by \obv{funs_is_function,funs}.
            Take $j\in\lambda$ such that $f(j)=\gamma$
                by \auto{fun_ran_iff}.
            Take $\delta\in\ran{p}$ such that $j\subseteq\delta$
                by \obv{cofinal_in}.
            $p$ is a function and $\dom{p}=\mu$
                by \obv{funs_is_function,funs}.
            Take $i\in\mu$ such that $p(i)=\delta$
                by \auto{fun_ran_iff}.
            Hence $p(i)\in\lambda$ by \obv{funs_type_apply}.
            $j$ and $p(i)$ are ordinals
                by \obv{limit_ordinal,elem_of_ordinal_is_ordinal}.
            Hence $j\in p(i)$ or $j=p(i)$
                by \obv{ordinal_subseteq_implies_elem_or_eq}.
            Hence $j\in\suc{p(i)}$ by \auto{suc_iff}.
            Hence $f(j)\in h(p(i))$
                by \auto{cofinal_envelope_dominates_prefix}.
            $f(j)$ and $h(p(i))$ are ordinals
                by \auto{funs_type_apply,limit_ordinal,elem_of_ordinal_is_ordinal}.
            Hence $f(j)\subseteq h(p(i))$
                by \obv{ordinal_elem_implies_subseteq}.
            $g(i)=h(p(i))$ by \auto{circ_apply}.
            Hence $\beta\subseteq g(i)$
                by \obv{subseteq_transitive}.
            $g$ is a function and $\dom{g}=\mu$
                by \auto{funs_is_function,funs}.
            Hence $g(i)\in\ran{g}$ by \auto{fun_ran_iff}.
        \end{subproof}
        Follows by \obv{cofinal_in}.
    \end{subproof}
    Hence $\cofinalMap{g}{\mu}{\alpha}$ by \obv{cofinal_map}.
    Hence $\cofinality{\alpha}\subseteq\card{\mu}$
        by \obv{cofinality_le_domain_card}.
    $\mu$ is a cardinal by \auto{cofinality_is_cardinal}.
    Hence $\card{\mu}=\mu$ by \obv{card_of_cardinal}.
    Hence $\lambda\subseteq\mu$.
    $\lambda$ is an ordinal by \auto{limit_ordinal}.
    $\mu=\cofinality{\lambda}$.
    $\mu\subseteq\lambda$ by \auto{cofinality_subseteq_ordinal}.
    Hence $\mu=\lambda$ by \obv{subseteq_antisymmetric}.
    Follows by \auto.
\end{proof}

\begin{definition}\label{regular_cardinal}
    $\kappa$ is a regular cardinal iff $\kappa$ is an infinite cardinal and
        $\cofinality{\kappa}=\kappa$.
\end{definition}

\begin{definition}\label{singular_cardinal}
    $\kappa$ is singular iff $\kappa$ is an infinite cardinal and
        $\kappa$ is not a regular cardinal.
\end{definition}

Here regular and singular cardinals are, by convention, infinite.  Thus the
adjective ``singular'' itself entails that its subject is an infinite cardinal;
finite cardinals lie outside this regular--singular terminology.

\begin{abbreviation}\label{uncountable_cofinality}
    $\kappa$ is of uncountable cofinality iff
        $\naturals\in\cofinality{\kappa}$.
\end{abbreviation}

\begin{theorem}\label{cofinality_of_limit_is_regular}
    Let $\alpha$ be a limit ordinal.
    Then $\cofinality{\alpha}$ is a regular cardinal.
\end{theorem}
\begin{proof}
    $\cofinality{\alpha}$ is an infinite cardinal
        by \obv{cofinality_of_limit_is_infinite_cardinal}.
    $\cofinality{\cofinality{\alpha}}=\cofinality{\alpha}$
        by \obv{cofinality_of_cofinality}.
    Follows by \obv{regular_cardinal}.
\end{proof}
